\section{Low-Dimensional Categories}

\begin{frame}
  \frametitle{Idea}

  \pause Introducing dimension fixes the issue.

  \pause Insight from Sophie's theory of higher dimensional structures:
  \begin{tightitemize}
    \item<+(1)-> Dimension (sizing) is intrinsic.
    \item<+(1)-> The formation rule of identity types decreases (size and)
      dimension.
  \end{tightitemize}

  \pause The formation rule of identity types should look like
  \begin{tightcenter}
    \begin{prooftree}
      \hypo{Γ\,⊢\,A\quad\text{$n$-type}}
      \hypo{Γ\,⊢\,x\,:\,A}
      \hypo{Γ\,⊢\,y\,:\,A}
      \infer3{Γ\,⊢\,x\,=\,y\quad\text{$(n-1)$-type}}
    \end{prooftree}
  \end{tightcenter}

  \pause In this work the relevant dimensions are -1 and 0 --- propositions and
  sets.
\end{frame}

\newcommand{\dimdeclentry}[2]{
    \AgdDta{#1}\AgdaSpace{}
  & \AgdaSymbol{:}\AgdaSpace{}
  & #2
}

\begin{frame}
  \frametitle{Dimension}

  \pause For our type theory with definitional proof-irrelevance, we will
  only need two dimensions: one for sets, and one for propositions.

  \pause Sets are $0$-dimensional, therefore the identity types of points of
  sets are $(-1)$-dimensional. We denote dimension $0$ and $-1$ as follows

  \pause \fbox{\Dim\AgdaSpace{}\AgdaSymbol{:}\AgdaSpace{}\AgdaPrimitive{Set}}
  \begin{tightcenter}
    \begin{prooftree}
      \infer0{\text{\IsDim{\DimNeg}}}
    \end{prooftree}
    \hspace{5em}
    \begin{prooftree}
      \infer0{\text{\IsDim{\DimZro}}}
    \end{prooftree}
  \end{tightcenter}

  \pause The well-scoped, untyped syntax now has signature
  \begin{tightcenter}
    \begin{tabular}{@{}l@{}l@{}l@{}}
         \dimdeclentry{\Ctx}{\AgdSet}
      \\ \dimdeclentry{\Var}{\AgdArr{\Ctx}{\AgdArr{\Dim}{\AgdSet}}}
      \\ \dimdeclentry{\Typ}{\AgdArr{\Ctx}{\AgdArr{\Dim}{\AgdSet}}}
      \\ \dimdeclentry{\Trm}{\AgdArr{\Ctx}{\AgdArr{\Dim}{\AgdSet}}}
    \end{tabular}
  \end{tightcenter}
\end{frame}

\begin{frame}
  \frametitle{Sets and Propositions}

  \begin{tightcenter}
    \begin{tabular}{lll}
         \pause
         \implicit{\IsTyp{\variable{P}}{\variable{Γ}}[\DimNeg]}
      &  \TypTyping{\variable{Γ}}{\variable{P}}
      &  \variable{P} is a proposition
      \\ \pause
         \implicit{\IsTyp{\variable{X}}{\variable{Γ}}[\DimZro]}
      &  \TypTyping{\variable{Γ}}{\variable{X}}
      &  \variable{X} is a set
      \\ \pause
         \implicit{\IsTrm{\variable{p}}{\variable{Γ}}[\DimNeg]}
      &  \TrmTyping{\variable{Γ}}{\variable{p}}{\variable{P}}
      &  \variable{p} is a proof of \variable{P}
      \\ \pause
         \implicit{\IsTrm{\variable{x}}{\variable{Γ}}[\DimZro]}
      &  \TrmTyping{\variable{Γ}}{\variable{x}}{\variable{P}}
      &  \variable{x} is a point of \variable{X}
    \end{tabular}
  \end{tightcenter}

  \pause The statement of proof-irrelevance for propositions is now
  \begin{tightcenter}
    \begin{prooftree}
      \hypo{
        \begin{tabular}{c}
             \text{\implicit{\IsCtx{\variable{Γ}}}}
          \\ \text{\CtxTyping{\variable{Γ}}}
        \end{tabular}
      }
      \hypo{
        \begin{tabular}{c}
             \text{\implicit{\IsTyp{\variable{P}}{\variable{Γ}}[\DimNeg]}}
          \\ \text{\TypTyping{\variable{Γ}}{\variable{P}}}
        \end{tabular}
      }
      \hypo{
        \begin{tabular}{c}
             \text{\implicit{\IsTrm{\variable{p}}{\variable{Γ}}[\DimNeg]}}
          \\ \text{\text{\TrmTyping{\variable{Γ}}{\variable{p}}{\variable{P}}}}
        \end{tabular}
      }
      \hypo{
        \begin{tabular}{c}
             \text{\implicit{\IsTrm{\variable{q}}{\variable{Γ}}[\DimNeg]}}
          \\ \text{\TrmTyping{\variable{Γ}}{\variable{q}}{\variable{P}}}
        \end{tabular}
      }
      \infer4{\text{\TrmConv{\variable{Γ}}{\variable{p}}{\variable{q}}{\variable{P}}}}
    \end{prooftree}
  \end{tightcenter}

  \pause And identity types for sets are now
  \begin{tightcenter}
    \begin{prooftree}
      \hypo{
        \begin{tabular}{c}
             \text{\implicit{\IsCtx{\variable{Γ}}}}
          \\ \text{\CtxTyping{\variable{Γ}}}
        \end{tabular}
      }
      \hypo{
        \begin{tabular}{c}
             \text{\implicit{\IsTyp{\variable{X}}{\variable{Γ}}[\DimZro]}}
          \\ \text{\TypTyping{\variable{Γ}}{\variable{X}}}
        \end{tabular}
      }
      \hypo{
        \begin{tabular}{c}
             \text{\implicit{\IsTrm{\variable{x}}{\variable{Γ}}[\DimZro]}}
          \\ \text{\text{\TrmTyping{\variable{Γ}}{\variable{x}}{\variable{X}}}}
        \end{tabular}
      }
      \hypo{
        \begin{tabular}{c}
             \text{\implicit{\IsTrm{\variable{y}}{\variable{Γ}}[\DimZro]}}
          \\ \text{\TrmTyping{\variable{Γ}}{\variable{y}}{\variable{X}}}
        \end{tabular}
      }
      %\infer4{\text{\TrmConv{\variable{Γ}}{\variable{t}}{\variable{u}}{\variable{A}}}}
      \infer4{\text{\TypTyping{\variable{Γ}}{\TypEq{\variable{x}}{\variable{y}}}}}
    \end{prooftree}
  \end{tightcenter}
\end{frame}

\begin{frame}
  \frametitle{Dimension-Generic Behavior}

  \pause More insight from Sophie's work: The formation rule of Π-types should
  look like
  \begin{tightcenter}
    \begin{prooftree}
      \hypo{Γ\,⊢\,A\quad\text{$n$-type}}
      \hypo{Γ,\,x\,:\,A\,⊢\,B\quad\text{$m$-type}}
      \infer2{Γ\,⊢\,Π\,x\,:\,A.\,B\quad\text{$m$-type}}
    \end{prooftree}
  \end{tightcenter}

  \pause The mechanized definition of the above is exactly the same:
  there is no reference to the specific value of the dimensions.

  \begin{tightcenter}
    \begin{prooftree}
      \hypo{\text{\implicit{\IsDim{\variable{d₀}\AgdaSpace{}\variable{d₁}}}}}
      \hypo{
        \begin{tabular}{c}
             \text{\implicit{\IsCtx{\variable{Γ}}}}
          \\ \text{\CtxTyping{\variable{Γ}}}
        \end{tabular}
      }
      \hypo{
        \begin{tabular}{c}
             \text{\implicit{\IsTyp{\variable{A}}{\variable{Γ}}[\variable{d₀}]}}
          \\ \text{\TypTyping{\variable{Γ}}{\variable{B}}}
        \end{tabular}
      }
      \hypo{
        \begin{tabular}{c}
             \text{\implicit{\IsTyp{\variable{B}}{\CtxExt{\variable{Γ}}{\variable{A}}}[\variable{d₁}]}}
          \\ \text{\TypTyping{\CtxExt{\variable{Γ}}{\variable{B}}}{\variable{B}}}
        \end{tabular}
      }
      \infer4{\text{\TypTyping{\variable{Γ}}{\TypArr{\variable{A}}{\variable{B}}}}}
    \end{prooftree}
  \end{tightcenter}
\end{frame}

